\bumppoints{26}

\question
The rabbitsfoot mussel was once widespread throughout rivers of eastern North America but now remains in only a few very small populations. The decrease occurred rapidly due to habitat loss. Genetics studies show that the remaining populations have little or no genetic variation. The best explanation for this result is:

\begin{choices}
\choice A bottleneck caused the loss of homozygosity. 
\choice Inbreeding caused an increase in heterozygosity. 
\choice Balancing selection favored only one or two alleles. 
\choice Directional selection favored only one or two alleles. 
\correctchoice Genetic drift caused the loss of heterozygosity. 
\end{choices}

\question
Retinitis pigmentosus (RP) is a rare form of blindness, but has a high frequency on Tristen de Cu\~nha, a small island in the South Atlantic Ocean.  The island was colonized by 15 Britons during the 1800s.  One person carried one allele for RP.  Today, of the 240 people descended from the original 15, 4 are blind and 9 are carriers for the disease.  Which of the following best explains the high proportion of RP on the island?

\begin{choices}
\choice bottleneck and inbreeding.
\choice founder effect and natural selection.
\choice genetic drift and mutation.
\choice inbreeding and mutation.
\correctchoice founder effect and inbreeding (a form of non-random mating).
\end{choices}

\question
The overuse of drugs such as penicillin has caused an increase in the number of bacterial species that are resistant to the drugs. The most likely explanation for this result is:

\begin{choices}
\choice Penicillin was not an effective drug. 
\choice The drugs caused mutations in the bacteria that gave them resistance. 
\choice Most bacteria can resist drugs. 
\choice Bacteria began making proteins resistant to the drugs. 
\correctchoice A few bacteria were naturally resistant and natural selection favored their increase. 
\end{choices}

\question
Which mouse has the greatest relative fitness? (A clutch is the offspring from one breeding event.)

\begin{choices}
\choice One that lives for 2 years, has 3 clutches per year, and 8 offspring per clutch. 
\choice One that lives for 4 years, has 3 clutches per year, and 4 offspring per clutch.
\choice One that lives for 3 years, has 3 clutches per year, and 5 offspring per clutch.
\correctchoice One that lives for 3 years, has 3 clutches per year, and 7 offspring per clutch. 
\choice One that lives for 3 years, has 2 clutches per year, and 6 offspring per clutch. 
\end{choices}

\question
According to Darwin's concept of descent with modification:

\begin{choices}
\correctchoice Species evolve from ancestral forms through the accumulation of different adaptations. 
\choice Populations evolve by natural selection acting on individuals. 
\choice All organisms overproduce offspring, leading to different adaptations. 
\choice Individuals with greater fitness will evolve more quickly than individuals with lower fitness. 
\choice All organisms overproduce offspring, leading to differences in survivorship. 
\end{choices}

\question
Over time, the movement of people among all countries on Earth has steadily increased. This has altered the course of human evolution by increasing

\begin{choices}
\choice natural selection. 
\correctchoice gene flow. 
\choice nonrandom mating. 
\choice genetic drift. 
\choice geographic isolation. 
\end{choices}

\question
Which \emph{one} of the following statements is true of natural selection?

\begin{choices}
\choice It results in descent with modification. 
\choice It requires genetic variation. 
\choice It involves differences in reproductive success. 
\correctchoice It requires genetic variation, results in descent with modification, and involves differences in reproductive success. 
\choice It results in descent with modification and involves differences in reproductive success. 
\end{choices}

\newpage

\question
\textit{Anopheles} mosquitoes, which carry the malaria parasite, cannot live above elevations of 5,900 feet. In addition, oxygen availability decreases with higher altitude. Consider a hypothetical human population that is adapted to life on the slopes of Mt. Kilimanjaro in Tanzania, a country in equatorial Africa. Mt. Kilimanjaro's base is about 2,600 feet above sea level and its peak is 19,341 feet above sea level. If the frequency of the sickle-cell allele in the population is plotted against altitude (feet above sea level), which of the following distributions is most likely, assuming little migration of people up or down the mountain?

\begin{multicols}{2}
	\begin{choices}
		\choice \includegraphics[width=0.35\textwidth]{anopA} 
		\choice \includegraphics[width=0.35\textwidth]{anopC} 
		\correctchoice \includegraphics[width=0.35\textwidth]{anopB} 
		\choice \includegraphics[width=0.35\textwidth]{anopD} 
	\end{choices}
\end{multicols}

\question
In a population with two alleles, $A_1$ and $A_2$, that is in Hardy-Weinberg equilibrium, the frequency of allele $A_1$ is 0.2. What is the frequency of $A_2A_2$ individuals? 

\begin{choices}
\choice 0.4 %dontshuffle
\choice 0.02 %dontshuffle
\choice 0.04 %dontshuffle
\choice 0.16 %dontshuffle
\choice 0.32 %dontshuffle
\correctchoice 0.64 %dontshuffle
\end{choices}

\question
Violation of which two Hardy-Weinberg assumptions \emph{always} causes overall homozygosity to increase in a population?

\begin{choices}
\correctchoice Genetic drift and non-random mating.
\choice Mutation and natural selection.
\choice Genetic drift and natural selection.
\choice Natural selection and gene flow.
\choice Gene flow and mutation.
\choice Non-random mating and gene flow.
\end{choices}

\question
Which mode of selection best applies to this example: female widowbirds choose males with the longest tails for mating. However, males with very long tails are caught more often by predators because they are slower to take flight. A long-term study showed that male tail length is not getting longer or shorter. 

\begin{choices}
\correctchoice stabilizing selection. 
\choice sexual selection. 
\choice disruptive selection. 
\choice balancing selection. 
\choice directional selection. 
\end{choices}

\question
Prairie chickens are endangered grassland birds. In Illinois, 
grasslands used to cover most of the state but farming and growth of 
cities changed the habitat over a fairly short period of time. Now, 
less than 1\% of the grasslands remain. A study showed that most genes 
had far fewer alleles than historic populations. The loss of genetic 
variation is most likely due to

\begin{choices}
\choice directional selection that favored survivorship in small habitats. 
\correctchoice a genetic bottleneck caused by the loss of habitat. 
\choice the small population size causing random fluctuation of allele frequencies. 
\choice a founder event caused by individuals moving to the remaining grasslands. 
\choice balancing selection, which maintains genetic variation in small populations. 
\end{choices}

\question
Which of these conditions should completely prevent the occurrence of natural selection in a population over time?

\begin{choices}
\choice The population lives in a habitat where there are no competing species present. 
\choice The environment is changing at a relatively slow rate. 
\choice The population size is large. 
\correctchoice All variation among individuals is due only to environmental factors. 
\choice The population mates non-randomly.
\end{choices}
